\section{Tiền xử lý dữ liệu}
\label{sec:preprocessing}

\subsection{Làm sạch dữ liệu bẩn (Noisy Data)}
\label{subsec:noisy_data}

Trong quá trình khám phá dữ liệu, nhóm phát hiện hai cột chứa dữ liệu bẩn nghiêm trọng do lỗi nhập liệu (\textit{data entry error}) hoặc hoán đổi cột (\textit{column swap}):

\begin{itemize}[noitemsep]
    \item \textbf{Sleep Duration}: Có 36 giá trị unique thay vì 4 giá trị hợp lệ. Khoảng 79 mẫu chứa các giá trị bất thường như tên thành phố (\texttt{Pune}), tên cột (\texttt{Work\_Study\_Hours}), hoặc số giờ vô nghĩa (\texttt{45-48 hours}).
    \item \textbf{Dietary Habits}: Có 23 giá trị unique thay vì 3 giá trị hợp lệ. Khoảng 27 mẫu chứa tên bằng cấp (\texttt{BSc}), giới tính (\texttt{Male}), hoặc các giá trị bị swap từ cột khác.
\end{itemize}

Chiến lược xử lý: ánh xạ (\textit{map}) các giá trị bẩn về category hợp lệ gần nhất dựa trên ngữ nghĩa; các giá trị không thể ánh xạ được thay thế bằng giá trị mode. Bảng~\ref{tab:noisy_data} tóm tắt kết quả trước và sau khi làm sạch.

\begin{table}[H]
\centering
\caption{Kết quả làm sạch dữ liệu bẩn}
\label{tab:noisy_data}
\begin{tabular}{lccl}
\toprule
\textbf{Cột} & \textbf{Unique (trước)} & \textbf{Unique (sau)} & \textbf{Ví dụ giá trị bẩn} \\
\midrule
Sleep Duration & 36 & 4 & Pune, 45-48 hours, 10-6 hours \\
Dietary Habits & 23 & 3 & BSc, Male, 1.0 \\
\bottomrule
\end{tabular}
\end{table}

\subsection{Xử lý Missing Values}
\label{subsec:missing_values}

Hình~\ref{fig:missing_values} trình bày tỷ lệ missing values theo từng cột. Ba cột có tỷ lệ missing cao nhất là Study Satisfaction, Academic Pressure và CGPA (đều 80,17\%), tiếp theo là Profession (26,03\%), Job Satisfaction và Work Pressure (đều 19,84\%).

\begin{figure}[H]
\centering
\includegraphics[width=0.85\textwidth]{preprocessing/missing_values.png}
\caption{Tỷ lệ Missing Values theo từng cột}
\label{fig:missing_values}
\end{figure}

Điểm quan trọng là các giá trị missing này \textbf{không phải ngẫu nhiên} mà thuộc loại \textbf{MNAR} (\textit{Missing Not At Random}). Hình~\ref{fig:missing_pattern} chứng minh rõ: Academic Pressure, CGPA, Study Satisfaction chỉ có giá trị ở nhóm Student (100\% missing ở Working Professional), trong khi Work Pressure và Job Satisfaction chỉ có ở Working Professional (100\% missing ở Student). Đây là thiết kế có chủ đích của bộ dữ liệu, không phải lỗi thu thập.

\begin{figure}[H]
\centering
\includegraphics[width=0.85\textwidth]{preprocessing/missing_pattern_by_group.png}
\caption{Pattern Missing Values theo nhóm Working Professional vs Student}
\label{fig:missing_pattern}
\end{figure}

\subsection{Gộp các cột điều kiện (Merge Conditional Columns)}
\label{subsec:merge_columns}

Do cơ chế MNAR được phát hiện ở Mục~\ref{subsec:missing_values}, việc điền missing bằng mean/median là \textbf{không phù hợp} --- vì giá trị missing mang ý nghĩa thực sự (sinh viên không có Work Pressure, người đi làm không có Academic Pressure). Nhóm áp dụng chiến lược gộp cột (\textit{merge}) như sau:

\begin{itemize}[noitemsep]
    \item \texttt{Academic Pressure} + \texttt{Work Pressure} $\rightarrow$ \texttt{Pressure}: gộp thành một biến thống nhất đo mức áp lực chung.
    \item \texttt{Study Satisfaction} + \texttt{Job Satisfaction} $\rightarrow$ \texttt{Satisfaction}: gộp thành biến đo mức hài lòng chung.
    \item \texttt{CGPA}: giữ nguyên, đồng thời tạo thêm biến nhị phân \texttt{has\_cgpa} để đánh dấu đối tượng là sinh viên hay không.
\end{itemize}

Kết quả: số lượng đặc trưng giảm từ 20 xuống 17, đồng thời loại bỏ hoàn toàn vấn đề missing values cho 5 cột điều kiện.

\subsection{Loại bỏ đặc trưng không cần thiết}
\label{subsec:drop_features}

Cột \texttt{Name} được loại bỏ vì chứa thông tin định danh cá nhân, không có giá trị dự đoán cho biến mục tiêu Depression.

\subsection{Mã hóa biến phân loại (Encoding Categorical Variables)}
\label{subsec:encoding}

Các biến phân loại được mã hóa theo hai phương pháp phù hợp:

\begin{itemize}[noitemsep]
    \item \textbf{Binary Encoding}: Áp dụng cho các biến nhị phân --- \texttt{Gender} (Male=1, Female=0), \texttt{Suicidal Thoughts} (Yes=1, No=0), \texttt{Family History} (Yes=1, No=0), \texttt{Working Professional or Student} (Student=0, Professional=1).
    \item \textbf{Ordinal Encoding}: Áp dụng cho các biến có thứ tự tự nhiên --- \texttt{Sleep Duration} (Less than 5h=0, 5-6h=1, 7-8h=2, More than 8h=3) và \texttt{Dietary Habits} (Unhealthy=0, Moderate=1, Healthy=2).
\end{itemize}

Lý do không sử dụng One-Hot Encoding: các biến này có thứ tự rõ ràng (ít $\rightarrow$ nhiều, xấu $\rightarrow$ tốt), ordinal encoding bảo toàn thông tin thứ tự này.

\subsection{Chuẩn hóa đặc trưng (Feature Scaling)}
\label{subsec:scaling}

Chuẩn hóa đặc trưng chỉ cần thiết cho các mô hình nhạy cảm với thang đo như Logistic Regression và SVM. Các mô hình tree-based (Decision Tree, Random Forest, XGBoost, LightGBM) không yêu cầu chuẩn hóa.

\textbf{Standardization} (Z-score normalization) --- dịch chuyển và co giãn dữ liệu sao cho trung bình bằng 0 và độ lệch chuẩn bằng 1. Lưu ý rằng phương pháp này \textbf{không} biến đổi hình dạng phân phối gốc về phân phối chuẩn (Gaussian), mà chỉ chuẩn hóa thang đo:

\begin{equation}
z = \frac{x - \mu}{\sigma}
\label{eq:standardization}
\end{equation}

trong đó $x$ là giá trị gốc, $\mu$ là trung bình mẫu và $\sigma$ là độ lệch chuẩn mẫu của đặc trưng.

\textbf{Min-Max Normalization} --- co giãn giá trị về khoảng $[0, 1]$:

\begin{equation}
x' = \frac{x - x_{\min}}{x_{\max} - x_{\min}}
\label{eq:minmax}
\end{equation}

trong đó $x_{\min}$ và $x_{\max}$ lần lượt là giá trị nhỏ nhất và lớn nhất của đặc trưng trong tập huấn luyện.

Trong dự án này, nhóm sử dụng \textbf{Standardization} (Phương trình~\ref{eq:standardization}) cho Logistic Regression và SVM vì phương pháp này ít bị ảnh hưởng bởi outliers hơn Min-Max.

\subsection{Kết quả tiền xử lý}
\label{subsec:preprocessing_result}

Bảng~\ref{tab:preprocessing_result} tóm tắt kết quả trước và sau quá trình tiền xử lý dữ liệu.

\begin{table}[H]
\centering
\caption{So sánh dữ liệu trước và sau tiền xử lý}
\label{tab:preprocessing_result}
\begin{tabular}{lcc}
\toprule
\textbf{Tiêu chí} & \textbf{Trước} & \textbf{Sau} \\
\midrule
Số đặc trưng & 20 & 17 \\
Số dòng (train) & 140.700 & 140.700 \\
Missing values & Có (5 cột, lên tới 80\%) & 0 \\
Noisy data & Có (2 cột) & 0 \\
Kiểu dữ liệu & Hỗn hợp (text + số) & Hoàn toàn số \\
\bottomrule
\end{tabular}
\end{table}
